%% Szkielet wygenerowany z docs/02_spis_tresci.md (plan v1, 11.11.2026).
%% Blok planroboczy usuń, gdy podrozdział jest napisany.

\chapter{Wyniki}
\label{ch:wyniki}

\begin{planroboczy}[ok. 12 s.]
Każdy podrozdział:
\begin{enumerate}
  \item pytanie badawcze,
  \item warunki (tabela albo zdanie),
  \item wynik (wykres i~tabela parametrów z~przedziałami ufności),
  \item krótka interpretacja bez odwołań do literatury (te trafiają do rozdziału 6).
\end{enumerate}

Liczby wyłącznie z \texttt{results/}.
\end{planroboczy}

\section{Weryfikacja narzędzia pomiarowego}
\label{sec:5-1}

\begin{planroboczy}[ok. 1 s.]
\begin{itemize}
  \item Powtarzalność przy stałym ziarnie.
  \item Odsetek odpowiedzi nieczytelnych, rozkład czasu wywołania.
  \item Porównanie z~trybem bez rozumowania (jeśli E0 to obejmie).
  \item Test odzyskiwania parametrów na danych syntetycznych.
\end{itemize}
\end{planroboczy}

\section{Progi bazowe: bateria i~wilgotność gleby}
\label{sec:5-2}

\begin{planroboczy}[ok. 2 s.]
\begin{itemize}
  \item PB1: krzywe psychometryczne dla promptu bazowego.
  \item Czy przełączenie jest progowe? Porównanie z~interpretacją jakościową z~pracy referencyjnej.
\end{itemize}
\end{planroboczy}

\section{Wpływ reprezentacji stanu}
\label{sec:5-3}

\begin{planroboczy}[ok. 1,5 s.]
\begin{itemize}
  \item PB4: etykiety vs liczby vs oba.
\end{itemize}
\end{planroboczy}

\section{Kalibracja zadanym progiem}
\label{sec:5-4}

\begin{planroboczy}[ok. 2 s.]
\begin{itemize}
  \item PB2: $\theta$ w~funkcji zadanego $X$. Nachylenie i~przesunięcie zależności.
\end{itemize}
\end{planroboczy}

\section{Ramowanie i~wskazówki zachowania}
\label{sec:5-5}

\begin{planroboczy}[ok. 2,5 s.]
\begin{itemize}
  \item PB3 i~PB5: przesunięcia $\theta$ i~zmiany $k$, $\gamma$, $\lambda$ względem promptu bazowego. Wykres „leśny” różnic z~przedziałami ufności.
\end{itemize}
\end{planroboczy}

\section{Rywalizacja bólów}
\label{sec:5-6}

\begin{planroboczy}[ok. 1,5 s.]
\begin{itemize}
  \item PB6: mapa prawdopodobieństwa wyboru bateria/gleba na siatce 2D. Granica decyzyjna dla 2–3 promptów.
\end{itemize}
\end{planroboczy}

\section{Walidacja w~pętli zamkniętej}
\label{sec:5-7}

\begin{planroboczy}[ok. 1,5 s.]
\begin{itemize}
  \item PB7: metryki epizodów naprawczych i~zależność od $\theta$ z~sondowania.
\end{itemize}
\end{planroboczy}
